\section{Кольца}

\begin{definition}
\emph{Кольцо}~--- это множество $R$ с двумя бинарными операциями $+$ и $\cdot$, для которых
\begin{enumerate}[label=(\alph*)]
\item $(R,+)$~--- коммутативная группа;
\item умножение ассоциативно, и существует элемент $1_R$ с $a\cdot1_R=a=1_R\cdot a$ для всех $a\in R$;
\item выполнены законы дистрибутивности: $(a+b)c=ac+bc$ и $a(b+c)=ab+ac$ для всех $a,b,c\in R$.
\end{enumerate}
\end{definition}

Обычно пишут $ab$ вместо $a\cdot b$ и $1$ вместо $1_R$. Если $1_R=0$, то $a=a\cdot1_R=a\cdot0=0$ для всех $a$, т.\,е.\
$R=\{0\}$. \emph{Подкольцо} кольца $R$~--- подмножество, содержащее $1_R$ и замкнутое относительно сложения, перехода к
противоположному элементу и умножения; оно само является кольцом. \emph{Гомоморфизм колец} $\alpha\colon R\to R'$~---
отображение с
\[
\alpha(a+b)=\alpha(a)+\alpha(b),\qquad \alpha(ab)=\alpha(a)\alpha(b),\qquad \alpha(1_R)=1_{R'}.
\]
Кольцо \emph{коммутативно}, если $ab=ba$ для всех $a,b$. Коммутативное кольцо $R$~--- \emph{область целостности}, если
$1_R\ne0$ и выполнен закон сокращения: из $ab=ac$ и $a\ne0$ следует $b=c$ (равносильно: из $ab=0$ следует $a=0$ или
$b=0$). \emph{Идеал} коммутативного кольца $R$~--- подгруппа $I\subset(R,+)$, для которой $ra\in I$ при $r\in R$, $a\in I$.
Идеал, порождённый элементами $a_1,\dots,a_n$, обозначается $(a_1,\dots,a_n)$; например, $(a)=aR$~---
\emph{главный идеал}.

Считаются известными поле частных области целостности, факторкольцо $R/I$ и однозначность разложения в евклидовых
кольцах. Факторкольцо $R/I$~--- область целостности
тогда и только тогда, когда идеал $I$ \emph{прост}, т.\,е.\ $I\ne R$ и из $ab\in I$ следует $a\in I$ или $b\in I$. В $\Z$
(и вообще в евклидовом кольце) каждый идеал $I\ne0$ порождается любым своим <<наименьшим>> ненулевым элементом, и
имеет место однозначное разложение на неприводимые множители. Через $\gcd(a,b)$ обозначается наибольший общий
делитель (в $\Z$~--- неотрицательный, в $F[X]$~--- унитарный); например, $\gcd(a,0)=|a|$ в $\Z$.

\begin{definition}[1.0a]
Пусть $R$~--- область целостности. Элемент $a\in R$ \emph{обратим}, если $ab=1$ для некоторого $b\in R$; элементы $a$ и
$ua$ с обратимым $u$ называются \emph{ассоциированными}. Элемент $a$ \emph{неприводим}, если он ненулевой, необратим и
допускает только тривиальные разложения: из $a=bc$ следует, что $b$ или $c$ обратим. Элемент $a$ \emph{прост}, если
$a\ne0$ и идеал $(a)$ прост, т.\,е.\ $a$ необратим и из $a\mid bc$ следует $a\mid b$ или $a\mid c$. Кольцо $R$
\emph{факториально}, если каждый ненулевой необратимый элемент раскладывается в произведение неприводимых, и это
разложение единственно с точностью до порядка множителей и замены их ассоциированными.
\end{definition}

\begin{proposition}[1.0b]
Пусть $R$~--- область целостности.
\begin{enumerate}[label=(\alph*)]
\item Каждый простой элемент $R$ неприводим.
\item Если $R$ факториально, то каждый неприводимый элемент $R$ прост. Итак, в факториальном кольце простые и
неприводимые элементы~--- одно и то же.
\end{enumerate}
\end{proposition}

\begin{proof}
(a) Пусть $a$ прост и $a=bc$. Тогда $bc\in(a)$, и, скажем, $b\in(a)$: $b=ak$. Получаем $a=akc$, и после сокращения на
$a\ne0$ имеем $kc=1$, т.\,е.\ $c$ обратим. Кроме того, $a\ne0$ и $a$ необратим (так как $(a)\ne R$). Значит, $a$ неприводим.

(b) Пусть $a$ неприводим. Он необратим, поэтому $(a)\ne R$. Пусть $bc\in(a)$, т.\,е.\ $bc=ak$. Если $b=0$ или $c=0$,
доказывать нечего. Иначе разложим $b$, $c$ и $k$ на неприводимые множители (обратимый элемент~--- пустое произведение с
обратимым множителем). Тогда $bc=ak$ даёт два разложения одного элемента; в одном из них участвует $a$, поэтому по
единственности $a$ ассоциирован с одним из неприводимых множителей $b$ или $c$. Значит, $a\mid b$ или $a\mid c$.
\end{proof}

\begin{remark}
Без факториальности утверждение (b) неверно. В кольце $\Z[\sqrt{-5}]=\{x+y\sqrt{-5}\mid x,y\in\Z\}$ имеем
$2\cdot3=(1+\sqrt{-5})(1-\sqrt{-5})$. Норма $N(x+y\sqrt{-5})=x^2+5y^2$ мультипликативна, и уравнение $x^2+5y^2=2$
неразрешимо в целых числах. Элемент $2$ необратим ($N(2)=4\ne1$, а у обратимого элемента $u$ из $uv=1$ следует
$N(u)N(v)=1$, т.\,е.\ $N(u)=1$), и в разложении $2=bc$ один из множителей имеет норму $1$, т.\,е.\ равен $\pm1$, так
что $2$ неприводим. Но $2$ не прост: он делит $(1+\sqrt{-5})(1-\sqrt{-5})=6$, но не делит ни один из множителей, так как
$\tfrac{1\pm\sqrt{-5}}2\notin\Z[\sqrt{-5}]$.
\end{remark}

\section{Поля}

\begin{definition}[1.1]
\emph{Поле}~--- это множество $F$ с бинарными операциями $+$ и $\cdot$, для которых
\begin{enumerate}[label=(\alph*)]
\item $(F,+)$~--- коммутативная группа;
\item $(F^\times,\cdot)$, где $F^\times=F\setminus\{0\}$,~--- коммутативная группа;
\item выполнен закон дистрибутивности.
\end{enumerate}
\end{definition}

Итак, поле~--- это ненулевое коммутативное кольцо, в котором каждый ненулевой элемент обратим. В частности, поле~---
область целостности: из $ab=0$, $a\ne0$ следует $b=a^{-1}ab=0$. В поле есть хотя бы два различных элемента $0$ и $1$;
наименьшее (и одно из самых важных) поле~--- $\F_2=\Z/2\Z=\{0,1\}$. \emph{Подполе}~--- подкольцо, замкнутое относительно
перехода к обратному элементу; оно само является полем.

\begin{lemma}[1.2]
Ненулевое коммутативное кольцо $R$ является полем тогда и только тогда, когда в нём нет идеалов, кроме $(0)$ и $R$.
\end{lemma}

\begin{proof}
Пусть $R$~--- поле и $I\ne0$~--- идеал. Если $0\ne a\in I$, то $1=a^{-1}a\in I$, и поэтому $I=R$. Обратно, пусть в $R$ нет
идеалов, кроме $(0)$ и $R$. Если $a\ne0$, то $(a)\ne(0)$, значит, $(a)=R\ni1$, т.\,е.\ $ab=1$ для некоторого $b\in R$.
\end{proof}

\begin{example}[1.3]
Полями являются $\Q$, $\R$, $\CC$ и $\F_p=\Z/p\Z$ ($p$ простое). Для $\F_p$: если $a\not\equiv0\pmod p$, то
$\gcd(a,p)=1$, и $ua+vp=1$ для некоторых $u,v\in\Z$; тогда $u$ (по модулю $p$)~--- обратный к $a$.
\end{example}

\emph{Гомоморфизм полей}~--- это просто гомоморфизм колец. Он всегда инъективен: его ядро~--- идеал, не содержащий $1$,
т.\,е.\ отличный от всего поля, и по \rf{1.2} он нулевой.

Пусть $F$~--- поле. \emph{$F$-алгебра} (алгебра над $F$)~--- кольцо $R$, содержащее $F$ в качестве подкольца.
\emph{Гомоморфизм $F$-алгебр} $\alpha\colon R\to R'$~--- гомоморфизм колец с $\alpha(c)=c$ для всех $c\in F$. $F$-алгебра $R$
\emph{конечна}, если она конечномерна как векторное пространство над $F$.

\section{Характеристика поля}

Пусть $F$~--- поле. Отображение
\[
\Z\to F,\qquad n\mapsto n\cdot1_F=\underbrace{1_F+1_F+\dots+1_F}_{n}\quad(n\ge0),\qquad (-n)\cdot1_F=-(n\cdot1_F),
\]
является гомоморфизмом колец. Например, $(m\cdot1_F)+(n\cdot1_F)=(m+n)\cdot1_F$ по ассоциативности сложения, а
$(m\cdot1_F)(n\cdot1_F)=mn\cdot1_F$ по дистрибутивности. Поэтому его ядро~--- идеал в $\Z$.

\emph{Случай 1: ядро равно $(0)$}, т.\,е.\ $n\cdot1_F=0\Rightarrow n=0$. Тогда ненулевые целые числа переходят в
обратимые элементы $F$, и отображение продолжается до гомоморфизма
\[
\Q\hookrightarrow F,\qquad \frac mn\mapsto(m\cdot1_F)(n\cdot1_F)^{-1}
\]
(корректность: если $m/n=m'/n'$, то $mn'=m'n$, и $(m\cdot1_F)(n\cdot1_F)^{-1}=(m'\cdot1_F)(n'\cdot1_F)^{-1}$). В этом случае
$F$ содержит копию $\Q$, и говорят, что $F$ имеет \emph{характеристику нуль}.

\emph{Случай 2: ядро не равно $(0)$}, т.\,е.\ $n\cdot1_F=0$ для некоторого $n\ne0$. Пусть $p$~--- наименьшее положительное
такое $n$; ядро~--- идеал в $\Z$, и он порождается числом $p$. Число $p$ простое: $p\ne1$, так как $1_F\ne0$, а если бы
$p=ab$ с $1<a,b<p$, то $(a\cdot1_F)(b\cdot1_F)=p\cdot1_F=0$, хотя оба множителя ненулевые по минимальности $p$,~---
противоречие с отсутствием делителей нуля. Итак, $n\mapsto n\cdot1_F$ задаёт изоморфизм $\Z/p\Z$ на подкольцо
$\{m\cdot1_F\mid m\in\Z\}$ поля $F$. В этом случае $F$ содержит копию $\F_p$, и говорят, что $F$ имеет
\emph{характеристику $p$}. Пишут $\operatorname{char}F=0$ или $\operatorname{char}F=p$.

Поле, изоморфное одному из полей $\F_2,\F_3,\F_5,\dots,\Q$, называется \emph{простым}. В простом поле нет собственных
подполей: подполе содержит $1$, а значит, и все $m\cdot1$, и (в случае $\Q$) все их частные. Поэтому \emph{каждое поле
содержит ровно одно простое подполе}: это образ $\F_p$ или $\Q$, построенный выше; он содержится в любом подполе, и
любое простое подполе, содержа его, с ним совпадает.

\begin{remark}[1.4]
Более общо, коммутативное кольцо $R$ \emph{имеет характеристику $p$} (соответственно $0$), если оно содержит простое поле
характеристики $p$ (соответственно $0$) в качестве подкольца. Тогда это простое поле единственно и по определению
содержит $1_R$. Характеристика есть не у каждого кольца: она есть ровно тогда, когда кольцо содержит какое-нибудь поле в
качестве подкольца. Например, её нет у $\Z$ (в $\Z$ нет подполей) и у $\F_2\times\F_3$: подполе содержало бы единицу
$(1,1)$, у которой $2\cdot(1,1)\ne0$, $3\cdot(1,1)\ne0$, но $6\cdot(1,1)=0$, что невозможно ни в характеристике $0$, ни в
характеристике $p$.

Пусть $R$~--- ненулевое коммутативное кольцо. Если $\operatorname{char}R=p\ne0$, то
\[
pa\overset{\text{def}}=\underbrace{a+\dots+a}_{p}=(\underbrace{1_R+\dots+1_R}_{p})\,a=0\cdot a=0
\quad\text{для всех }a\in R.
\]
Обратно, если $pa=0$ для всех $a\in R$, то ядро гомоморфизма $\Z\to R$, $n\mapsto n\cdot1_R$, содержит $p$ и не равно $\Z$
(так как $1_R\ne0$), т.\,е.\ равно $p\Z$; образ~--- подкольцо, изоморфное $\F_p$, и $\operatorname{char}R=p$.
\end{remark}

\subsection*{Эндоморфизм Фробениуса}

В ненулевом коммутативном кольце $R$ обычная индукция доказывает формулу бинома
\[
(a+b)^m=a^m+\binom m1a^{m-1}b+\binom m2a^{m-2}b^2+\dots+b^m.
\]
Если $p$ простое и $1\le r\le p-1$, то $p$ делит $\binom pr=\dfrac{p!}{r!\,(p-r)!}$: в равенстве
$p!=\binom pr\,r!\,(p-r)!$ простое $p$ делит левую часть, но не делит $r!\,(p-r)!$ (все сомножители меньше $p$),
и потому делит $\binom pr$. Следовательно, если $\operatorname{char}R=p$, то
\[
(a+b)^p=a^p+b^p\qquad\text{для всех }a,b\in R,
\]
и отображение $a\mapsto a^p\colon R\to R$ является гомоморфизмом колец (и даже $\F_p$-алгебр: на элементах
$m\cdot1_R$ оно тождественно, так как $(1+\dots+1)^p=1^p+\dots+1^p$; это малая теорема Ферма). Оно называется
\emph{эндоморфизмом Фробениуса} кольца $R$. Отображение $a\mapsto a^{p^n}$, $n\ge1$,~--- композиция $n$ эндоморфизмов
Фробениуса, поэтому тоже гомоморфизм. Следовательно,
\[
(a_1+\dots+a_m)^{p^n}=a_1^{p^n}+\dots+a_m^{p^n}\qquad\text{для всех }a_i\in R.
\]
Если $F$~--- поле, то эндоморфизм Фробениуса инъективен (как любой гомоморфизм полей; непосредственно: $a^p=0\Rightarrow
a=0$), и потому является автоморфизмом, если $F$ конечно (инъективное отображение конечного множества в себя
биективно). Для бесконечных полей это, вообще говоря, не так (пример \rf{1.4a}), хотя бывает и так: на любом
алгебраически замкнутом поле характеристики $p$ (скажем, на алгебраическом замыкании поля $\F_p$, \rf{1.44})
Фробениус сюръективен, так как у многочлена $X^p-a$ есть корень.

\begin{example}[1.4a]
В поле рациональных функций $\F_p(T)$ (см.\ \rf{1.10}) элемент $T$ не является $p$-й степенью, так что Фробениус не
сюръективен. Действительно, пусть $T=(f/g)^p$ с $f,g\in\F_p[T]$, $g\ne0$. Тогда $Tg^p=f^p$, и степени дают
$1+p\deg g=p\deg f$, что невозможно по модулю $p$.
\end{example}

\emph{Характеристическая экспонента} поля $F$ равна $1$, если $\operatorname{char}F=0$, и равна $p$, если
$\operatorname{char}F=p\ne0$. Итак, если $q$~--- характеристическая экспонента $F$ и $n\ge1$, то $x\mapsto x^{q^n}$~---
изоморфизм поля $F$ на подполе $F^{q^n}=\{x^{q^n}\mid x\in F\}$.

\section{Кольца многочленов}

Пусть $F$~--- поле.

\begin{definition}[1.5]
Кольцо $F[X]$ многочленов от символа (<<переменной>>) $X$ с коэффициентами в $F$~--- это векторное пространство над $F$
с базисом $1,X,\dots,X^n,\dots$ и умножением
\[
\Bigl(\sum_ia_iX^i\Bigr)\Bigl(\sum_jb_jX^j\Bigr)=\sum_k\Bigl(\sum_{i+j=k}a_ib_j\Bigr)X^k.
\]
\end{definition}

\emph{Универсальное свойство.} Пусть $R$~--- $F$-алгебра и $r\in R$~--- элемент, перестановочный со всеми элементами $F$
(например, $R$ коммутативна). Тогда существует единственный гомоморфизм $F$-алгебр $\alpha\colon F[X]\to R$ с
$\alpha(X)=r$, а именно $\alpha\bigl(\sum a_iX^i\bigr)=\sum a_ir^i$ (<<подстановка $r$>>; пишут $f(r)=\alpha(f)$).
Единственность ясна: гомоморфизм $F$-алгебр обязан переводить $\sum a_iX^i$ в $\sum a_ir^i$. Аддитивность этого отображения
очевидна, а мультипликативность следует из правила умножения, так как $a_ir^ib_jr^j=a_ib_jr^{i+j}$ (здесь нужна
перестановочность $r$ с $b_j$).

Для $f=a_nX^n+\dots+a_0$ с $a_n\ne0$ пишут $\deg f=n$ (\emph{степень}) и называют $a_n$ \emph{старшим коэффициентом};
$f$ \emph{унитарен}, если $a_n=1$. Удобно считать $\deg0=-\infty$. Старший коэффициент произведения равен произведению
старших коэффициентов (в поле нет делителей нуля), поэтому
\[
\deg(fg)=\deg f+\deg g,\qquad \deg(f+g)\le\max(\deg f,\deg g),
\]
кольцо $F[X]$~--- область целостности, а его обратимые элементы~--- ненулевые константы.

\begin{proposition}[1.6: деление с остатком]
Пусть $f,g\in F[X]$, $g\ne0$. Существуют $q,r\in F[X]$, для которых
\[
f=gq+r,\qquad r=0\ \text{или}\ \deg r<\deg g;
\]
более того, $q$ и $r$ определены однозначно.
\end{proposition}

\begin{proof}
\emph{Существование}~--- индукцией по $\deg f$. Если $\deg f<\deg g$, берём $q=0$, $r=f$. Пусть $\deg f=n\ge m=\deg g$,
$a$ и $b$~--- старшие коэффициенты $f$ и $g$. Многочлен $f_1=f-\frac ab X^{n-m}g$ имеет степень $<n$, и по индукции
$f_1=gq_1+r$ с $\deg r<\deg g$; тогда $f=g\bigl(q_1+\frac abX^{n-m}\bigr)+r$.

\emph{Единственность.} Если $gq+r=gq'+r'$, то $g(q-q')=r'-r$. Справа степень $<\deg g$, а слева при $q\ne q'$ степень
$\ge\deg g$. Значит, $q=q'$ и $r=r'$.
\end{proof}

Таким образом, $F[X]$~--- евклидово кольцо с нормой $\deg$, а потому кольцо главных идеалов и факториальное кольцо.
Многочлен из $F[X]$ называется \emph{неприводимым}, если он непостоянен и не является произведением двух многочленов
меньшей степени.

\begin{remark}[1.6a]
Это определение согласовано с общим (\rf{1.0a}): обратимые элементы $F[X]$~--- ненулевые константы, поэтому разложение
$f=gh$ тривиально (один из множителей обратим) ровно тогда, когда один из множителей имеет степень $0$, т.\,е.\ другой
имеет степень $\deg f$. По \rf{1.0b} неприводимые многочлены~--- это в точности простые элементы $F[X]$.
\end{remark}

\begin{remark}[1.7]
Пусть $f\in F[X]$ непостоянен и $a\in F$. Деление с остатком на $X-a$ даёт
\[
f=(X-a)q+c,\qquad q\in F[X],\ c\in F,
\]
и подстановка $X=a$ показывает, что $c=f(a)$. Поэтому, если $a$~--- \emph{корень} $f$ (т.\,е.\ $f(a)=0$), то $X-a$ делит $f$.
Отсюда следует, что \emph{у $f$ не более $\deg f$ корней} (ср.\ упражнение 1-3). Это можно вывести из однозначности
разложения, а можно прямо индукцией по $\deg f$ для всех ненулевых $f$ (у ненулевой константы корней нет): если
$a$~--- корень, то $f=(X-a)q$ с $\deg q=\deg f-1$, и для любого
другого корня $b\ne a$ из $0=f(b)=(b-a)q(b)$ получаем $q(b)=0$; по предположению индукции таких $b$ не больше
$\deg f-1$.

Здесь существенно, что $F$~--- поле (достаточно области целостности): в $\Z/8\Z$ у многочлена $X^2-1$ четыре корня
$1,3,5,7$.
\end{remark}

\subsection*{Производные и формула Тейлора}

Для $f=\sum_{i=0}^na_iX^i\in F[X]$ и целого $k\ge0$ положим
\[
D^{(k)}f=\sum_{i=k}^n\binom ika_iX^{i-k}\qquad(\text{так что }D^{(0)}f=f).
\]
Многочлен $f'=D^{(1)}f=\sum_{i\ge1}ia_iX^{i-1}$~--- \emph{(формальная) производная} $f$, а $D^{(k)}f$ называют
\emph{$k$-й производной Хассе}. Отображения $f\mapsto D^{(k)}f$ $F$-линейны. Через $f^{(k)}$ обозначим $k$-кратную
производную ($f^{(0)}=f$, $f^{(k+1)}=(f^{(k)})'$). Так как $(X^i)^{(k)}=i(i-1)\cdots(i-k+1)X^{i-k}=k!\binom ikX^{i-k}$
(при $k>i$ обе части равны нулю), по линейности
\[
f^{(k)}=k!\,D^{(k)}f .
\]
Всё сказанное дословно переносится на многочлены с коэффициентами в любом коммутативном кольце (биномиальные
коэффициенты~--- целые числа); например, для $f\in\Z[X]$ все $D^{(k)}f$ лежат в $\Z[X]$.

\begin{proposition}[1.7a: формула Тейлора]
Пусть $f\in F[X]$, $\deg f\le n$, а $x,h$~--- элементы коммутативной $F$-алгебры $R$. Тогда
\[
f(x+h)=\sum_{k=0}^n\bigl(D^{(k)}f\bigr)(x)\,h^k=f(x)+f'(x)\,h+\bigl(D^{(2)}f\bigr)(x)\,h^2+\dots
\]
В частности, для $a\in F$ в $F[X]$
\[
f(X+a)=\sum_{k=0}^n\bigl(D^{(k)}f\bigr)(a)\,X^k,\qquad f(X)=\sum_{k=0}^n\bigl(D^{(k)}f\bigr)(a)\,(X-a)^k,
\]
а если $\operatorname{char}F=0$, то $D^{(k)}f=f^{(k)}/k!$ и
\[
f(X+a)=f(a)+f'(a)X+\frac{f''(a)}{2!}X^2+\dots+\frac{f^{(n)}(a)}{n!}X^n .
\]
\end{proposition}

\begin{proof}
В коммутативном кольце $R$ по формуле бинома $(x+h)^i=\sum_{k=0}^i\binom ikx^{i-k}h^k$, поэтому
\[
f(x+h)=\sum_{i=0}^na_i\sum_{k=0}^i\binom ikx^{i-k}h^k=\sum_{k=0}^n\Bigl(\sum_{i=k}^n\binom ika_ix^{i-k}\Bigr)h^k
=\sum_{k=0}^n\bigl(D^{(k)}f\bigr)(x)\,h^k .
\]
Два частных случая получаются при $R=F[X]$, $x=a$ и $h=X$, соответственно $h=X-a$ (тогда $x+h=X$). Если
$\operatorname{char}F=0$, то $k!\cdot1_F\ne0$, и из $f^{(k)}=k!\,D^{(k)}f$ получаем $D^{(k)}f=f^{(k)}/k!$.
\end{proof}

Например, коэффициент при $h^2$ равен $\sum_i\binom i2a_ix^{i-2}$, и так как $\binom i2=\frac{i(i-1)}2$, это
$\frac12\sum_ii(i-1)a_ix^{i-2}=\frac{f''(x)}{2}$; так же $\binom ik=\frac{i(i-1)\cdots(i-k+1)}{k!}$ объясняет
появление $k!$ в знаменателе. В характеристике $p$ запись через $f^{(k)}/k!$ теряет смысл при $k\ge p$ (там $k!=0$),
и обычных производных не хватает: для $f=X^p$ все $f^{(k)}$ с $k\ge1$ равны нулю, тогда как $(X+a)^p=X^p+a^p$; здесь
$D^{(p)}f=1$. Производные Хассе работают в любой характеристике.

Формула Тейлора в $\Z[X]$ даёт, например, для $f\in\Z[X]$, простого $p$ и целых $a$, $s$, $k\ge1$:
\[
f(a+p^ks)\equiv f(a)+p^ks\,f'(a)\pmod{p^{2k}},
\]
так как остальные слагаемые $\bigl(D^{(j)}f\bigr)(a)\,p^{jk}s^j$, $j\ge2$, делятся на $p^{2k}$. На этом основан подъём
корней по модулю $p^k$ (ср.\ \rf{1.18e} и упражнение 1-8).

\begin{theorem}[1.8: алгоритм Евклида]
Пусть $f,g\in F[X]$ ненулевые и хотя бы один из них непостоянен; $d=\gcd(f,g)$~--- их унитарный наибольший общий
делитель. Тогда существуют $a,b\in F[X]$, для которых
\[
a f+b g=d,\qquad \deg a<\deg g,\qquad \deg b<\deg f.
\]
\end{theorem}

Напомним, как работает алгоритм. Можно считать, что $\deg f\ge\deg g$ (иначе меняем роли $f$ и $g$). Положим
$r_0=f$, $r_1=g$ и, пока остаток не обратится в нуль, делим с остатком:
\[
\begin{aligned}
r_0&=q_1r_1+r_2,\\
r_1&=q_2r_2+r_3,\\
&\ \,\vdots\\
r_{n-2}&=q_{n-1}r_{n-1}+r_n,\\
r_{n-1}&=q_nr_n,
\end{aligned}
\qquad \deg r_1>\deg r_2>\dots>\deg r_n\ge0,
\]
где $r_n$~--- последний ненулевой остаток (степени остатков строго убывают, поэтому процесс обрывается; положим
$r_{n+1}=0$). Тогда $r_n$ делит $r_{n-1}$, а значит, и $r_{n-2}=q_{n-1}r_{n-1}+r_n$, и так далее вверх, вплоть до $r_1=g$ и
$r_0=f$.
С другой стороны, двигаясь снизу вверх,
\[
r_n=r_{n-2}-q_{n-1}r_{n-1}=r_{n-2}-q_{n-1}(r_{n-3}-q_{n-2}r_{n-2})=\dots=a_0f+b_0g,
\]
так что каждый общий делитель $f$ и $g$ делит $r_n$. Итак, $r_n$~--- наибольший общий делитель, и
$d=r_n/\ell$, где $\ell$~--- старший коэффициент $r_n$.

\begin{proof}[Доказательство оценок степеней]
Пусть $a_0f+b_0g=d$ (например, полученное выше соотношение, делённое на $\ell$). Разделим $a_0$ на $g$ с остатком:
$a_0=gq+a$, $\deg a<\deg g$ (если $g$~--- константа, то $a=0$), и положим $b=b_0+qf$. Тогда
\[
af+bg=(a_0-gq)f+(b_0+qf)g=a_0f+b_0g=d .
\]
Оценим $\deg b$. Имеем $bg=d-af$, и
\[
\deg b+\deg g=\deg(d-af)\le\max(\deg d,\ \deg a+\deg f).
\]
Оба числа под максимумом меньше $\deg f+\deg g$: первое~--- потому что $\deg d\le\min(\deg f,\deg g)$, а один из
многочленов непостоянен; второе~--- потому что $\deg a<\deg g$. Значит, $\deg b<\deg f$.
\end{proof}

Условие <<хотя бы один непостоянен>> нужно: если $f,g$~--- ненулевые константы, то $d=1$, а неравенства требуют
$a=b=0$.

\subsection*{Матричная запись алгоритма Евклида}

Оценки степеней можно получить и прямо из алгоритма, причём точные. Шаг $r_{i-1}=q_ir_i+r_{i+1}$ записывается
матрично:
\[
\begin{pmatrix}r_i\\ r_{i+1}\end{pmatrix}=N_i\begin{pmatrix}r_{i-1}\\ r_i\end{pmatrix},\qquad
N_i=\begin{pmatrix}0&1\\1&-q_i\end{pmatrix},\qquad \det N_i=-1,\qquad
N_i^{-1}=M_i=\begin{pmatrix}q_i&1\\1&0\end{pmatrix}.
\]
Положим $s_0=1$, $t_0=0$, $s_1=0$, $t_1=1$ и
\[
s_{i+1}=s_{i-1}-q_is_i,\qquad t_{i+1}=t_{i-1}-q_it_i\qquad(1\le i\le n).
\]

\begin{proposition}[1.8a]
Для $0\le k\le n$
\[
N_kN_{k-1}\cdots N_1=\begin{pmatrix}s_k&t_k\\ s_{k+1}&t_{k+1}\end{pmatrix},\qquad\text{и потому}\qquad
s_kf+t_kg=r_k\quad(0\le k\le n+1).
\]
В частности,
\[
\begin{pmatrix}r_n\\0\end{pmatrix}=N_n\cdots N_1\begin{pmatrix}f\\ g\end{pmatrix},\qquad
\begin{pmatrix}f\\ g\end{pmatrix}=M_1\cdots M_n\begin{pmatrix}r_n\\0\end{pmatrix}.
\]
Отсюда снова видно, что $r_n$ делит $f$ и $g$ (второе равенство: у матрицы $M_1\cdots M_n$ полиномиальные элементы) и
что $r_n=s_nf+t_ng$ (первое равенство).
\end{proposition}

\begin{proof}
При $k=0$ слева пустое произведение $E$, а справа $\left(\begin{smallmatrix}s_0&t_0\\s_1&t_1\end{smallmatrix}\right)=E$.
Шаг индукции:
\[
N_{k+1}\begin{pmatrix}s_k&t_k\\ s_{k+1}&t_{k+1}\end{pmatrix}
=\begin{pmatrix}s_{k+1}&t_{k+1}\\ s_k-q_{k+1}s_{k+1}&t_k-q_{k+1}t_{k+1}\end{pmatrix}
=\begin{pmatrix}s_{k+1}&t_{k+1}\\ s_{k+2}&t_{k+2}\end{pmatrix}.
\]
Применяя эти матрицы к столбцу $\binom fg=\binom{r_0}{r_1}$, получаем $\binom{r_k}{r_{k+1}}$, т.\,е.\ $s_kf+t_kg=r_k$. Второе
выделенное равенство получается умножением первого на $(N_n\cdots N_1)^{-1}=M_1\cdots M_n$.
\end{proof}

Обычный <<обратный ход>> (подстановка остатков снизу вверх) вычисляет то же произведение матриц, только слева направо:
$r_n=(0\ \ 1)\,N_{n-1}\cdots N_1\binom fg$. По ассоциативности результат тот же, но при <<прямом ходе>> (справа
налево, т.\,е.\ рекуррентное вычисление $s_i,t_i$ вместе с остатками) выражения не нужно хранить и раскрывать.

Обозначим $n_i=\deg r_i$; тогда $n_0\ge n_1>n_2>\dots>n_n\ge0$ и $\deg q_i=n_{i-1}-n_i$, так что $\deg q_i\ge1$ при
$i\ge2$ (лишь $q_1$ может быть константой, если $\deg f=\deg g$).

\begin{lemma}[1.8b]
$\deg t_i=n_0-n_{i-1}$ при $1\le i\le n+1$ и $\deg s_i=n_1-n_{i-1}$ при $2\le i\le n+1$.
\end{lemma}

\begin{proof}
Ключевое наблюдение: в рекуррентных формулах старшие члены никогда не сокращаются.

\emph{Для $t$.} $t_1=1$ имеет степень $0=n_0-n_0$; $t_2=-q_1$ имеет степень $n_0-n_1$. Пусть $2\le i\le n$ и формула
верна для $t_{i-1}$, $t_i$. Тогда $\deg(q_it_i)=(n_{i-1}-n_i)+(n_0-n_{i-1})=n_0-n_i$, а
$\deg t_{i-1}=n_0-n_{i-2}<n_0-n_i$, так как $n_{i-2}\ge n_{i-1}>n_i$. Поэтому $\deg t_{i+1}=\deg(t_{i-1}-q_it_i)=n_0-n_i$.

\emph{Для $s$.} $s_2=s_0-q_1s_1=1$ имеет степень $0=n_1-n_1$, и (если $n\ge2$) $s_3=s_1-q_2s_2=-q_2$ имеет степень
$n_1-n_2$. При $3\le i\le n$: $\deg(q_is_i)=(n_{i-1}-n_i)+(n_1-n_{i-1})=n_1-n_i$, а $\deg s_{i-1}=n_1-n_{i-2}<n_1-n_i$,
так как $i-2\ge1$ и $n_{i-2}>n_i$. Поэтому $\deg s_{i+1}=n_1-n_i$.
\end{proof}

\begin{corollary}[1.8c]
Пусть $n\ge2$ (т.\,е.\ $g\nmid f$). Тогда многочлены $a=s_n/\ell$, $b=t_n/\ell$ удовлетворяют равенству $af+bg=d$ и
\[
\deg a=n_1-n_{n-1}<\deg g-\deg d,\qquad \deg b=n_0-n_{n-1}<\deg f-\deg d .
\]
Если же $n=1$ (т.\,е.\ $g\mid f$), то $d=g/\ell$, $a=0$, $b=1/\ell$.
\end{corollary}

\begin{proof}
Равенство $af+bg=d$~--- это \rf{1.8a}, делённое на $\ell$. Неравенства следуют из \rf{1.8b} при $i=n$, так как
$n_{n-1}>n_n=\deg d$, $n_1=\deg g$, $n_0=\deg f$.
\end{proof}

Таким образом, алгоритм Евклида сам выдаёт коэффициенты Безу с оценками даже более сильными, чем в \rf{1.8}:
<<не хватает>> ровно $n_{n-1}-n_n=\deg q_n\ge1$.

\begin{example}[1.8d]
Пусть $f=X^4+2X^2+1=(X^2+1)^2$, $g=X^3+X+1$ над $\Q$. Алгоритм Евклида:
\[
\begin{aligned}
X^4+2X^2+1&=X\cdot(X^3+X+1)+(X^2-X+1), & q_1&=X,\\
X^3+X+1&=(X+1)(X^2-X+1)+X, & q_2&=X+1,\\
X^2-X+1&=(X-1)\cdot X+1, & q_3&=X-1,\\
X&=X\cdot1, & q_4&=X.
\end{aligned}
\]
Итак, $n=4$, $r_4=1=d$, степени остатков $(n_0,\dots,n_4)=(4,3,2,1,0)$. Прямой ход:
\begin{center}
\begin{tabular}{@{}cllll@{}}
\toprule
$i$ & $r_i$ & $q_i$ & $s_i$ & $t_i$\\
\midrule
0 & $X^4+2X^2+1$ & & $1$ & $0$\\
1 & $X^3+X+1$ & $X$ & $0$ & $1$\\
2 & $X^2-X+1$ & $X+1$ & $1$ & $-X$\\
3 & $X$ & $X-1$ & $-(X+1)$ & $X^2+X+1$\\
4 & $1$ & $X$ & $X^2$ & $-(X^3+X-1)$\\
5 & $0$ & & $-(X^3+X+1)$ & $X^4+2X^2+1$\\
\bottomrule
\end{tabular}
\end{center}
Например, $s_4=s_2-q_3s_3=1+(X-1)(X+1)=X^2$. Степени $s_2,\dots,s_5$ равны $0,1,2,3=n_1-n_{i-1}$, степени
$t_1,\dots,t_5$ равны $0,1,2,3,4=n_0-n_{i-1}$, как в \rf{1.8b}. Получаем
\[
X^2\cdot(X^4+2X^2+1)-(X^3+X-1)(X^3+X+1)=1,
\]
т.\,е.\ $a=X^2$, $b=-(X^3+X-1)$, $\deg a=2<3=\deg g$, $\deg b=3<4=\deg f$. Проверка: $(X^3+X-1)(X^3+X+1)=(X^3+X)^2-1$. Обратный ход даёт то же самое:
\[
1=r_2-(X-1)r_3=r_2-(X-1)\bigl(r_1-(X+1)r_2\bigr)=X^2r_2-(X-1)r_1=X^2r_0-(X^3+X-1)r_1 .
\]
\end{example}

\begin{aside}[1.8e: континуанты]
Произведение матриц $M_i$ описывается явно. \emph{Континуанты}~--- многочлены от коммутирующих переменных, заданные
рекуррентно:
\[
K_{-1}=0,\qquad K_0=1,\qquad K_k(x_1,\dots,x_k)=x_kK_{k-1}(x_1,\dots,x_{k-1})+K_{k-2}(x_1,\dots,x_{k-2}).
\]
Так, $K_1=x_1$, $K_2=x_1x_2+1$, $K_3=x_1x_2x_3+x_1+x_3$, $K_4=x_1x_2x_3x_4+x_1x_2+x_1x_4+x_3x_4+1$.

\emph{Правило Эйлера:} $K_k(x_1,\dots,x_k)$~--- сумма всех мономов, получающихся из $x_1x_2\cdots x_k$ вычёркиванием
любого набора непересекающихся пар \emph{соседних} множителей $x_jx_{j+1}$ (включая пустой набор). Действительно, обозначим
эту сумму $S_k$; мономы, в которых $x_k$ не вычеркнут, дают $x_kS_{k-1}$, а мономы, в которых вычеркнута пара
$x_{k-1}x_k$, дают $S_{k-2}$; итак, $S_k$ удовлетворяет той же рекурренте с теми же начальными значениями. В частности,
$K_k$ не меняется при обращении порядка переменных, и
\[
\begin{pmatrix}x_1&1\\1&0\end{pmatrix}\cdots\begin{pmatrix}x_k&1\\1&0\end{pmatrix}
=\begin{pmatrix}K_k(x_1,\dots,x_k)&K_{k-1}(x_1,\dots,x_{k-1})\\ K_{k-1}(x_2,\dots,x_k)&K_{k-2}(x_2,\dots,x_{k-1})\end{pmatrix}
\]
(индукция по $k$: умножение справа на $\left(\begin{smallmatrix}x_{k+1}&1\\1&0\end{smallmatrix}\right)$ переводит столбцы
$(u,v)$ в $(x_{k+1}u+v,\ u)$, а это и есть рекуррента). Определитель левой части равен $(-1)^k$. Обращая эту матрицу при
$x_i=q_i$, $k=n$, и сравнивая с \rf{1.8a}, получаем
\begin{gather*}
f=K_n(q_1,\dots,q_n)\,r_n,\qquad g=K_{n-1}(q_2,\dots,q_n)\,r_n,\\
s_n=(-1)^nK_{n-2}(q_2,\dots,q_{n-1}),\qquad t_n=(-1)^{n+1}K_{n-1}(q_1,\dots,q_{n-1}).
\end{gather*}
Если $\deg x_j\ge1$ при $j\ge2$, то по правилу Эйлера старший член $K_k(x_1,\dots,x_k)$~--- это старший член
$x_1\cdots x_k$ (вычёркивание любой пары понижает степень). Отсюда снова $\deg s_n=\sum_{i=2}^{n-1}\deg q_i=n_1-n_{n-1}$ и
$\deg t_n=\sum_{i=1}^{n-1}\deg q_i=n_0-n_{n-1}$. В примере \rf{1.8d}: $s_4=K_2(q_2,q_3)=(X+1)(X-1)+1=X^2$ и
$t_4=-K_3(q_1,q_2,q_3)=-(X(X+1)(X-1)+X+X-1)=-(X^3+X-1)$.
\end{aside}

\begin{remark}[1.9]
Пусть $I$~--- ненулевой идеал в $F[X]$ и $f$~--- ненулевой многочлен наименьшей степени в $I$. Тогда $I=(f)$: для
$h\in I$ запишем $h=fq+r$ с $\deg r<\deg f$; тогда $r=h-fq\in I$, и по минимальности $r=0$. Если выбрать $f$ унитарным,
то он определён идеалом $I$ однозначно (два унитарных образующих делят друг друга и потому совпадают). Итак, ненулевые
идеалы $F[X]$ находятся во взаимно однозначном соответствии с унитарными многочленами. Ненулевые простые идеалы
соответствуют унитарным неприводимым многочленам: $(f)$ прост $\iff$ $f$~--- простой элемент $\iff$ $f$ неприводим
(\rf{1.0b}, \rf{1.6a}). (Нулевой идеал тоже прост, так как $F[X]$~--- область целостности.)
\end{remark}

\begin{remark}[1.10]
Так как $F[X]$~--- область целостности, можно образовать её поле частных $F(X)$~--- \emph{поле рациональных функций}.
Его элементы~--- частные $f/g$ многочленов, $g\ne0$, причём $f/g=f'/g'$ тогда и только тогда, когда $fg'=f'g$.
\end{remark}
